\documentclass[11pt,a4paper]{article}

% ===== PACKAGES =====
\usepackage[utf8]{inputenc}
\usepackage[T1]{fontenc}
\usepackage[margin=1in]{geometry}
\usepackage{hyperref}
\usepackage{listings}
\usepackage{xcolor}
\usepackage{booktabs}
\usepackage{array}
\usepackage{longtable}
\usepackage{fancyhdr}
\usepackage{titlesec}
\usepackage{tcolorbox}
\usepackage{fontawesome5}
\usepackage{amssymb}
\usepackage{amsmath}
\usepackage{enumitem}
\usepackage{graphicx}

% ===== COLORS =====
\definecolor{codegreen}{rgb}{0.25,0.49,0.48}
\definecolor{codegray}{rgb}{0.5,0.5,0.5}
\definecolor{codepurple}{rgb}{0.58,0,0.82}
\definecolor{codeblue}{rgb}{0.13,0.29,0.53}
\definecolor{codeorange}{rgb}{0.8,0.4,0.0}
\definecolor{backcolour}{rgb}{0.95,0.95,0.95}
\definecolor{warningbg}{rgb}{1.0,0.95,0.85}
\definecolor{warningborder}{rgb}{0.9,0.6,0.2}
\definecolor{tipbg}{rgb}{0.9,0.97,0.9}
\definecolor{tipborder}{rgb}{0.3,0.7,0.3}
\definecolor{notebg}{rgb}{0.9,0.93,1.0}
\definecolor{noteborder}{rgb}{0.3,0.4,0.7}

% ===== IEC 61131-3 CODE LISTING STYLE =====
\lstdefinelanguage{IEC61131}{
    keywords={VAR, END_VAR, VAR_INPUT, VAR_OUTPUT, VAR_IN_OUT, VAR_GLOBAL, 
              TYPE, END_TYPE, STRUCT, END_STRUCT, FUNCTION, END_FUNCTION,
              FUNCTION_BLOCK, END_FUNCTION_BLOCK, PROGRAM, END_PROGRAM,
              IF, THEN, ELSE, ELSIF, END_IF, CASE, OF, END_CASE,
              FOR, TO, BY, DO, END_FOR, WHILE, END_WHILE, REPEAT, UNTIL, END_REPEAT,
              TRUE, FALSE, AND, OR, NOT, XOR, MOD, RETURN,
              ARRAY, OF, POINTER, REFERENCE, TO, AT, EXTENDS, IMPLEMENTS,
              METHOD, END_METHOD, PROPERTY, END_PROPERTY,
              INTERFACE, END_INTERFACE,
              PUBLIC, PRIVATE, PROTECTED, INTERNAL, FINAL, ABSTRACT,
              THIS, SUPER, VAR_STAT, CONSTANT, VAR_INST},
    keywordstyle=\color{codeblue}\bfseries,
    keywords=[2]{BOOL, BYTE, WORD, DWORD, LWORD, SINT, USINT, INT, UINT, 
                 DINT, UDINT, LINT, ULINT, REAL, LREAL, STRING, WSTRING,
                 TIME, LTIME, DATE, TIME_OF_DAY, TOD, DATE_AND_TIME, DT,
                 TON, TOF, TP, CTU, CTD, CTUD, R_TRIG, F_TRIG,
                 CONCAT, CONCAT2, LEFT, RIGHT, MID, LEN, FIND, INSERT, DELETE, REPLACE},
    keywordstyle=[2]\color{codepurple}\bfseries,
    keywords=[3]{ADR, SIZEOF, REF, ABS, SQRT, SIN, COS, TAN, LN, LOG, EXP,
                 SEL, MUX, LIMIT, MAX, MIN, MOVE, ADSLOGSTR},
    keywordstyle=[3]\color{codeorange}\bfseries,
    sensitive=false,
    morecomment=[l]{//},
    morecomment=[s]{(*}{*)},
    commentstyle=\color{codegreen}\itshape,
    stringstyle=\color{codeorange},
    morestring=[b]',
    morestring=[b]",
}

\lstdefinestyle{iecstyle}{
    language=IEC61131,
    backgroundcolor=\color{backcolour},
    basicstyle=\ttfamily\small,
    breakatwhitespace=false,
    breaklines=true,
    captionpos=b,
    keepspaces=true,
    numbers=left,
    numbersep=8pt,
    numberstyle=\tiny\color{codegray},
    showspaces=false,
    showstringspaces=false,
    showtabs=false,
    tabsize=4,
    frame=single,
    framerule=0.5pt,
    rulecolor=\color{codegray},
    xleftmargin=15pt,
    framexleftmargin=15pt,
}

\lstset{style=iecstyle}

% ===== TCOLORBOX STYLES =====
\tcbuselibrary{skins,breakable}

\newtcolorbox{warningbox}{
    colback=warningbg,
    colframe=warningborder,
    boxrule=1pt,
    left=10pt,
    right=10pt,
    top=8pt,
    bottom=8pt,
    breakable,
    before upper={\faExclamationTriangle\ \textbf{Warning:} }
}

\newtcolorbox{tipbox}{
    colback=tipbg,
    colframe=tipborder,
    boxrule=1pt,
    left=10pt,
    right=10pt,
    top=8pt,
    bottom=8pt,
    breakable,
    before upper={\faLightbulb\ \textbf{Tip:} }
}

\newtcolorbox{notebox}{
    colback=notebg,
    colframe=noteborder,
    boxrule=1pt,
    left=10pt,
    right=10pt,
    top=8pt,
    bottom=8pt,
    breakable,
    before upper={\faInfoCircle\ \textbf{Note:} }
}

% ===== HEADER/FOOTER =====
\pagestyle{fancy}
\fancyhf{}
\fancyhead[L]{\leftmark}
\fancyhead[R]{TwinCAT 3 Lecture Notes}
\fancyfoot[C]{\thepage}
\renewcommand{\headrulewidth}{0.4pt}
\renewcommand{\footrulewidth}{0.4pt}

% ===== HYPERREF SETUP =====
\hypersetup{
    colorlinks=true,
    linkcolor=codeblue,
    filecolor=magenta,
    urlcolor=cyan,
    pdftitle={TwinCAT 3 - Standard Functions and Function Blocks},
    pdfauthor={},
    bookmarks=true,
}

% ===== TITLE FORMATTING =====
\titleformat{\section}
    {\normalfont\Large\bfseries}{\thesection}{1em}{}[{\titlerule[0.8pt]}]
\titleformat{\subsection}
    {\normalfont\large\bfseries}{\thesubsection}{1em}{}
\titleformat{\subsubsection}
    {\normalfont\normalsize\bfseries}{\thesubsubsection}{1em}{}

% ===== DOCUMENT INFO =====
\title{
    \vspace{-1cm}
    \textbf{Lecture Notes: TwinCAT 3}\\
    \Large Standard Functions and Function Blocks (Tc2\_Standard)
}
\author{}
\date{\today}

% ===== BEGIN DOCUMENT =====
\begin{document}

\maketitle

\tableofcontents
\newpage

% ============================================================
\section{Introduction to the IEC 61131-3 Standard}
% ============================================================

This lecture focuses on the basic functions and function blocks defined by the \textbf{IEC 61131-3 standard}, which are essential for every PLC programmer. In TwinCAT 3, these are found in the \textbf{Tc2\_Standard library}. Because these functions are standardized, code written using them is \textbf{highly portable} across different PLC environments.

\begin{table}[h]
\centering
\begin{tabular}{@{}lll@{}}
\toprule
\textbf{Category} & \textbf{Function Blocks} & \textbf{Purpose} \\
\midrule
Timers & TON, TOF, TP & Time-based control \\
Counters & CTU, CTD, CTUD & Event counting \\
Edge Detection & R\_TRIG, F\_TRIG & Signal transition detection \\
String Operations & CONCAT, LEFT, RIGHT, MID & Text manipulation \\
\bottomrule
\end{tabular}
\caption{Key IEC 61131-3 Standard Library Components}
\end{table}

\begin{notebox}
The Tc2\_Standard library is automatically referenced in most new TwinCAT projects. If you encounter ``undefined symbol'' errors for timers or triggers, verify that the library reference exists under References in your PLC project.
\end{notebox}

% ============================================================
\section{Timer Function Blocks}
% ============================================================

Timers are used to manage time-sensitive operations, such as waiting for a process to complete or limiting the duration of an action.

\subsection{Timer On Delay (TON)}

\begin{itemize}[leftmargin=*]
    \item \textbf{Purpose:} The \texttt{TON} function block turns its output (\texttt{Q}) on only after a specified time delay (\texttt{PT}).
    
    \item \textbf{Inputs:}
    \begin{itemize}
        \item \textbf{IN:} A rising edge (transition from FALSE to TRUE) starts the timer. A falling edge resets the timer.
        \item \textbf{PT (Preset Time):} The duration of the delay (e.g., \texttt{T\#3s}).
    \end{itemize}
    
    \item \textbf{Outputs:}
    \begin{itemize}
        \item \textbf{Q:} Becomes TRUE once the elapsed time reaches the preset time.
        \item \textbf{ET (Elapsed Time):} Shows how much time has passed since the timer started.
    \end{itemize}
    
    \item \textbf{Behavior:} If the input \texttt{IN} goes low before the preset time is reached, the timer resets and \texttt{Q} never goes high.
\end{itemize}

\subsubsection*{Example: TON Declaration and Call}

\begin{lstlisting}
// PROGRAM: MAIN
// Demonstrates TON (Timer On Delay) function block usage
// The timer output Q goes TRUE after the preset time elapses

PROGRAM MAIN
VAR
    // =====================================================
    // TON Instance Declaration
    // The timer can be initialized with a preset time (PT)
    // This sets the default delay duration
    // =====================================================
    fbTimer     : TON := (PT := T#10S);  // 10 second delay timer
    
    // Alternative declaration without initialization
    fbTimer2    : TON;                    // PT must be set when called
    
    // Additional timers for different purposes
    fbStartupDelay  : TON := (PT := T#5S);   // 5 second startup delay
    fbTimeoutTimer  : TON := (PT := T#30S);  // 30 second timeout
    
    // Control variables
    bStartProcess   : BOOL;     // Input: Start button
    bProcessActive  : BOOL;     // Output: Process is running
    bTimerComplete  : BOOL;     // Output: Timer has finished
    
    // Monitoring variables
    tElapsedTime    : TIME;     // Current elapsed time
END_VAR

// =====================================================
// Basic TON Call
// The timer starts counting when IN transitions to TRUE
// Q becomes TRUE after PT time has elapsed
// =====================================================

// Call the timer with the start button as trigger
fbTimer(IN := bStartProcess);

// Read the outputs
bTimerComplete := fbTimer.Q;      // TRUE when 10 seconds elapsed
tElapsedTime := fbTimer.ET;       // Current elapsed time (0 to PT)

// =====================================================
// Alternative: Set PT dynamically when calling
// =====================================================
fbTimer2(
    IN := bStartProcess,
    PT := T#15S                   // Override preset time
);

// =====================================================
// Practical Example: Startup Delay
// Wait 5 seconds after start button before activating
// =====================================================
fbStartupDelay(IN := bStartProcess);

IF fbStartupDelay.Q THEN
    bProcessActive := TRUE;       // Activate after delay
END_IF

// =====================================================
// Using ET for progress indication
// ET counts from T#0S up to PT while IN is TRUE
// =====================================================
// tElapsedTime can be used for:
// - Progress bars on HMI
// - Logging how long a process takes
// - Conditional logic based on partial timing
\end{lstlisting}

\begin{lstlisting}
// Extended example: Multiple timers working together

VAR
    fbWarmupTimer   : TON := (PT := T#2M);    // 2 minute warmup
    fbCycleTimer    : TON := (PT := T#500MS); // 500ms cycle
    fbAlarmDelay    : TON := (PT := T#5S);    // 5 second alarm delay
    
    bMachineOn      : BOOL;
    bOverTemp       : BOOL;
    bWarmupComplete : BOOL;
    bAlarmActive    : BOOL;
END_VAR

// Warmup timer - machine needs 2 minutes before operation
fbWarmupTimer(IN := bMachineOn);
bWarmupComplete := fbWarmupTimer.Q;

// Alarm delay - only trigger alarm if overtemp persists for 5 seconds
// This prevents false alarms from brief temperature spikes
fbAlarmDelay(IN := bOverTemp AND bWarmupComplete);
bAlarmActive := fbAlarmDelay.Q;
\end{lstlisting}

\begin{figure}[h]
\centering
\begin{tabular}{|c|c|c|c|c|c|c|c|c|}
\hline
\textbf{Time} & 0s & 1s & 2s & 3s & 4s & 5s & 6s & 7s \\
\hline
\textbf{IN} & 0 & 1 & 1 & 1 & 1 & 1 & 0 & 0 \\
\hline
\textbf{ET} & 0 & 1s & 2s & 3s & 3s & 3s & 0 & 0 \\
\hline
\textbf{Q} (PT=3s) & 0 & 0 & 0 & 1 & 1 & 1 & 0 & 0 \\
\hline
\end{tabular}
\caption{TON Timing Diagram (PT = 3 seconds)}
\end{figure}

\subsection{Timer Off Delay (TOF)}

\begin{itemize}[leftmargin=*]
    \item \textbf{Purpose:} The inverse of \texttt{TON}. It counts time when the input \texttt{IN} is \textbf{low}.
    \item \textbf{Trigger:} It is triggered by a \textbf{falling edge}.
    \item \textbf{Output Q:} Goes TRUE immediately when \texttt{IN} is high and only drops to FALSE after the input has been low for the duration of \texttt{PT}.
\end{itemize}

\begin{table}[h]
\centering
\begin{tabular}{@{}llll@{}}
\toprule
\textbf{Timer} & \textbf{Starts When} & \textbf{Q Goes HIGH} & \textbf{Q Goes LOW} \\
\midrule
TON & IN rises & After PT elapsed & When IN falls \\
TOF & IN falls & When IN rises & After PT elapsed \\
TP & IN rises & When IN rises & After PT elapsed \\
\bottomrule
\end{tabular}
\caption{Comparison of Timer Function Blocks}
\end{table}

\subsection{Use Case: Motor Safety}

Timers can be used as a \textbf{safety feature} to prevent equipment from running indefinitely due to system errors. For instance, if an HMI button controlling a motor gets ``stuck'' in a high state (due to a crash or communication error), a \texttt{TON} can ensure the motor moves for a maximum of 5 or 10 seconds before automatically stopping.

\subsubsection*{Example: Motor Safety Logic}

\begin{lstlisting}
// PROGRAM: MotorSafetyControl
// Demonstrates using a timer as a safety limit for motor operation
// Prevents motor from running indefinitely if HMI crashes

PROGRAM MotorSafetyControl
VAR
    // =====================================================
    // Safety Timer
    // Limits maximum motor run time to prevent damage
    // =====================================================
    fbSafetyTimer   : TON := (PT := T#10S);  // 10 second maximum run time
    
    // Inputs
    bOperatorClick  : BOOL;     // HMI button: Operator commands motor to move
    bLimitSwitch    : BOOL;     // Physical limit switch (motor reached position)
    bSafetyOK       : BOOL := TRUE;  // Safety circuit status
    
    // Outputs
    bMoveMotor      : BOOL;     // Output to motor contactor/drive
    bTimeoutAlarm   : BOOL;     // Alarm: Motor ran too long
    
    // Status
    bMotorMoving    : BOOL;     // Motor is currently moving
END_VAR

// =====================================================
// Start the safety timer when operator clicks the button
// The timer runs as long as the operator holds the button
// =====================================================
fbSafetyTimer(IN := bOperatorClick);

// =====================================================
// MOTOR SAFETY LOGIC
// Motor runs when:
//   - Operator is pressing the button (bOperatorClick = TRUE)
//   - AND safety timer has NOT elapsed (fbSafetyTimer.Q = FALSE)
//   - AND safety circuit is OK
//   - AND limit switch not reached
// =====================================================
bMoveMotor := bOperatorClick 
              AND NOT fbSafetyTimer.Q 
              AND bSafetyOK 
              AND NOT bLimitSwitch;

// Alternative simplified version:
// bMoveMotor := (bOperatorClick AND NOT fbSafetyTimer.Q);

// =====================================================
// Generate alarm if timer elapsed while button still pressed
// This indicates the HMI may be stuck or motor is jammed
// =====================================================
bTimeoutAlarm := fbSafetyTimer.Q AND bOperatorClick;

// Status output
bMotorMoving := bMoveMotor;

// =====================================================
// BEHAVIOR EXPLANATION:
// 1. Operator presses button -> Motor starts, timer starts
// 2. If operator releases before 10s -> Motor stops, timer resets
// 3. If 10 seconds elapse with button held -> Motor stops automatically
// 4. Operator must release and re-press button to restart
// =====================================================
\end{lstlisting}

\begin{lstlisting}
// Extended example with multiple safety features

VAR
    fbMaxRunTimer   : TON := (PT := T#30S);   // Maximum run time
    fbJogTimer      : TON := (PT := T#2S);    // Jog mode limit
    
    bJogForward     : BOOL;
    bJogReverse     : BOOL;
    bAutoMode       : BOOL;
    bMotorFwd       : BOOL;
    bMotorRev       : BOOL;
END_VAR

// Jog mode - limited to 2 seconds per press
fbJogTimer(IN := bJogForward OR bJogReverse);

// Motor outputs with safety limits
bMotorFwd := bJogForward AND NOT fbJogTimer.Q AND NOT bJogReverse;
bMotorRev := bJogReverse AND NOT fbJogTimer.Q AND NOT bJogForward;
\end{lstlisting}

\begin{warningbox}
Always implement safety timers for manual motor control. If an HMI freezes or loses communication while a button is pressed, the motor could run indefinitely, causing equipment damage, product loss, or personnel injury.
\end{warningbox}

% ============================================================
\section{Edge Detection Function Blocks}
% ============================================================

Edge detectors are used to identify the exact moment a signal changes state.

\subsection{Rising Edge (R\_TRIG)}

\begin{itemize}[leftmargin=*]
    \item \textbf{Purpose:} Detects a transition from FALSE to TRUE.
    \item \textbf{Behavior:} The output \texttt{Q} stays TRUE for only \textbf{one PLC cycle}. To re-trigger it, the input must return to FALSE and then go TRUE again.
    \item \textbf{Common Use:} Detecting an operator button press or raising an alarm only once when a condition is met.
\end{itemize}

\subsection{Falling Edge (F\_TRIG)}

\begin{itemize}[leftmargin=*]
    \item \textbf{Purpose:} Detects a transition from TRUE to FALSE.
    \item \textbf{Behavior:} Like \texttt{R\_TRIG}, the output \texttt{Q} is high for exactly one PLC cycle upon detecting the falling edge.
\end{itemize}

\subsubsection*{Example: Edge Trigger Counter}

\begin{lstlisting}
// PROGRAM: EdgeDetectionDemo
// Demonstrates R_TRIG and F_TRIG for edge detection
// Essential for counting events and preventing repeated actions

PROGRAM EdgeDetectionDemo
VAR
    // =====================================================
    // Edge Detection Function Block Instances
    // Each instance maintains its own state (previous CLK value)
    // =====================================================
    rtButtonPress   : R_TRIG;    // Rising edge detector for button
    ftButtonRelease : F_TRIG;    // Falling edge detector
    rtSensorTrigger : R_TRIG;    // Rising edge for sensor
    
    // Inputs
    bButton         : BOOL;      // Pushbutton input
    bSensorSignal   : BOOL;      // Sensor input (e.g., part detected)
    
    // Counters
    nTriggerCounter : INT := 0;  // Counts rising edges (button presses)
    nReleaseCounter : INT := 0;  // Counts falling edges (button releases)
    nPartCount      : INT := 0;  // Counts parts detected by sensor
    
    // Status flags
    bJustPressed    : BOOL;      // TRUE for one cycle when pressed
    bJustReleased   : BOOL;      // TRUE for one cycle when released
END_VAR

// =====================================================
// Call the rising edge detector
// CLK is the signal to monitor for edges
// =====================================================
rtButtonPress(CLK := bButton);

// =====================================================
// Use the Q output to increment a counter
// Q is TRUE for exactly ONE PLC cycle on each rising edge
// This ensures the counter increments once per button press
// =====================================================
IF rtButtonPress.Q THEN
    // This block executes exactly ONCE per button press
    nTriggerCounter := nTriggerCounter + 1;
    bJustPressed := TRUE;
ELSE
    bJustPressed := FALSE;
END_IF

// =====================================================
// Falling edge detection - detect button release
// =====================================================
ftButtonRelease(CLK := bButton);

IF ftButtonRelease.Q THEN
    // This block executes exactly ONCE per button release
    nReleaseCounter := nReleaseCounter + 1;
    bJustReleased := TRUE;
ELSE
    bJustReleased := FALSE;
END_IF

// =====================================================
// Practical Example: Part Counter
// Count parts as they pass a sensor
// =====================================================
rtSensorTrigger(CLK := bSensorSignal);

IF rtSensorTrigger.Q THEN
    nPartCount := nPartCount + 1;
    // Could also log the event, update HMI, etc.
END_IF
\end{lstlisting}

\begin{lstlisting}
// Extended example: Multiple edge detectors for different purposes

VAR
    rtStartButton   : R_TRIG;
    rtStopButton    : R_TRIG;
    rtResetButton   : R_TRIG;
    ftSafetyGate    : F_TRIG;    // Detect when gate opens (signal drops)
    
    bStart          : BOOL;
    bStop           : BOOL;
    bReset          : BOOL;
    bGateClosed     : BOOL;
    
    bMachineRunning : BOOL;
    nCycleCount     : UDINT := 0;
    bGateAlarm      : BOOL;
END_VAR

// Start/Stop logic using edge detection
rtStartButton(CLK := bStart);
rtStopButton(CLK := bStop);
rtResetButton(CLK := bReset);
ftSafetyGate(CLK := bGateClosed);

// Start machine on start button press
IF rtStartButton.Q AND bGateClosed THEN
    bMachineRunning := TRUE;
END_IF

// Stop machine on stop button press
IF rtStopButton.Q THEN
    bMachineRunning := FALSE;
END_IF

// Alarm if gate opens while machine is running
IF ftSafetyGate.Q AND bMachineRunning THEN
    bGateAlarm := TRUE;
    bMachineRunning := FALSE;  // Emergency stop
END_IF

// Reset alarm
IF rtResetButton.Q THEN
    bGateAlarm := FALSE;
    nCycleCount := 0;
END_IF
\end{lstlisting}

\begin{figure}[h]
\centering
\begin{tabular}{|c|c|c|c|c|c|c|c|c|c|}
\hline
\textbf{Cycle} & 1 & 2 & 3 & 4 & 5 & 6 & 7 & 8 & 9 \\
\hline
\textbf{CLK} & 0 & 1 & 1 & 1 & 0 & 0 & 1 & 1 & 0 \\
\hline
\textbf{R\_TRIG.Q} & 0 & \textbf{1} & 0 & 0 & 0 & 0 & \textbf{1} & 0 & 0 \\
\hline
\textbf{F\_TRIG.Q} & 0 & 0 & 0 & 0 & \textbf{1} & 0 & 0 & 0 & \textbf{1} \\
\hline
\end{tabular}
\caption{Edge Detection Timing Diagram}
\end{figure}

\subsection{Use Case: Preventing Message Flooding}

When logging events (e.g., using \texttt{ADSLOGSTR}), calling the function directly with a boolean input can result in hundreds of messages per second because the PLC runs cyclically. Using an \texttt{R\_TRIG} ensures that only \textbf{one message is generated} per event (like a button click or temperature threshold), and the signal must be reset before another message can be sent.

\begin{lstlisting}
// Example: Using R_TRIG to prevent log message flooding

VAR
    rtLogTrigger    : R_TRIG;
    bAlarmCondition : BOOL;     // TRUE when temperature too high
END_VAR

rtLogTrigger(CLK := bAlarmCondition);

IF rtLogTrigger.Q THEN
    // This logs ONCE when alarm condition becomes true
    // Not every cycle while it remains true
    ADSLOGSTR(ADSLOG_MSGTYPE_WARN, 'Temperature alarm activated', '');
END_IF
\end{lstlisting}

\begin{tipbox}
Always use R\_TRIG or F\_TRIG when you need to perform an action ``once per event'' rather than ``continuously while condition is true.'' This is especially important for logging, counting, and state machine transitions.
\end{tipbox}

% ============================================================
\section{String Manipulation: CONCAT}
% ============================================================

The \texttt{CONCAT} function is used to join two strings together into one.

\begin{itemize}[leftmargin=*]
    \item \textbf{Limitations:} The standard \texttt{CONCAT} function works with strings up to \textbf{255 characters}.
    \item \textbf{Arbitrary Length:} For strings longer than 255 characters, the \texttt{CONCAT2} function should be used instead.
\end{itemize}

\subsubsection*{Example: String Concatenation}

\begin{lstlisting}
// PROGRAM: StringOperationsDemo
// Demonstrates CONCAT and other string manipulation functions
// Useful for building messages, file paths, and formatted output

PROGRAM StringOperationsDemo
VAR
    // =====================================================
    // String Variables for Concatenation
    // STRING default length is 80 characters
    // STRING(n) specifies maximum length of n characters
    // =====================================================
    s1          : STRING := 'Hello ';       // First string
    s2          : STRING := 'World!';       // Second string
    sResult     : STRING(255);              // Result string (larger buffer)
    
    // Additional strings for examples
    sPrefix     : STRING := 'Error: ';
    sMessage    : STRING := 'Sensor disconnected';
    sFullMessage: STRING(255);
    
    // For building file paths
    sBasePath   : STRING := 'C:\Data\';
    sFileName   : STRING := 'log_';
    sExtension  : STRING := '.txt';
    sFullPath   : STRING(255);
    
    // For dynamic messages
    sProductID  : STRING := 'WIDGET-001';
    sStatus     : STRING := ' - PASSED';
    sLogEntry   : STRING(255);
    
    // Numeric values to include in strings
    nErrorCode  : INT := 42;
    sErrorCode  : STRING;
END_VAR

// =====================================================
// Basic CONCAT Usage
// Joins two strings into one
// Syntax: result := CONCAT(string1, string2);
// =====================================================
sResult := CONCAT(s1, s2);
// sResult now contains: 'Hello World!'

// =====================================================
// Building an error message
// =====================================================
sFullMessage := CONCAT(sPrefix, sMessage);
// sFullMessage now contains: 'Error: Sensor disconnected'

// =====================================================
// Multiple concatenations (chained)
// CONCAT only takes 2 arguments, so chain calls for more
// =====================================================
sFullPath := CONCAT(sBasePath, sFileName);
sFullPath := CONCAT(sFullPath, sExtension);
// sFullPath now contains: 'C:\Data\log_.txt'

// Alternative: Nested CONCAT calls
sFullPath := CONCAT(CONCAT(sBasePath, sFileName), sExtension);

// =====================================================
// Building a log entry
// =====================================================
sLogEntry := CONCAT(sProductID, sStatus);
// sLogEntry now contains: 'WIDGET-001 - PASSED'

// =====================================================
// IMPORTANT: Converting numbers to strings
// Use INT_TO_STRING, REAL_TO_STRING, etc.
// =====================================================
sErrorCode := INT_TO_STRING(nErrorCode);
sFullMessage := CONCAT('Error code: ', sErrorCode);
// sFullMessage now contains: 'Error code: 42'
\end{lstlisting}

\begin{lstlisting}
// Extended example: Building formatted messages

VAR
    sTimestamp      : STRING := '2024-01-15 10:30:00';
    sMachineName    : STRING := 'Conveyor_01';
    sEventType      : STRING := 'FAULT';
    sDescription    : STRING := 'Motor overtemperature';
    
    sLogLine        : STRING(255);
    sTemp           : STRING(255);
END_VAR

// Build a structured log message:
// "[2024-01-15 10:30:00] Conveyor_01 - FAULT: Motor overtemperature"

sLogLine := CONCAT('[', sTimestamp);
sLogLine := CONCAT(sLogLine, '] ');
sLogLine := CONCAT(sLogLine, sMachineName);
sLogLine := CONCAT(sLogLine, ' - ');
sLogLine := CONCAT(sLogLine, sEventType);
sLogLine := CONCAT(sLogLine, ': ');
sLogLine := CONCAT(sLogLine, sDescription);

// Result: '[2024-01-15 10:30:00] Conveyor_01 - FAULT: Motor overtemperature'
\end{lstlisting}

\begin{table}[h]
\centering
\begin{tabular}{@{}llp{6cm}@{}}
\toprule
\textbf{Function} & \textbf{Syntax} & \textbf{Description} \\
\midrule
CONCAT & CONCAT(s1, s2) & Joins two strings \\
LEFT & LEFT(str, n) & Returns first n characters \\
RIGHT & RIGHT(str, n) & Returns last n characters \\
MID & MID(str, pos, n) & Returns n chars starting at pos \\
LEN & LEN(str) & Returns string length \\
FIND & FIND(str1, str2) & Finds str2 in str1, returns position \\
\bottomrule
\end{tabular}
\caption{Common String Manipulation Functions}
\end{table}

\begin{notebox}
When concatenating strings, always ensure your result variable has sufficient capacity. If the combined string exceeds the variable's declared length, it will be truncated without warning.
\end{notebox}

% ============================================================
\section{Library Management}
% ============================================================

To access these functions, your TwinCAT project must have a reference to the \textbf{Tc2\_Standard library}. Most new TwinCAT projects include this reference automatically, but it is necessary for functions like timers and triggers to be recognized by the compiler.

\begin{table}[h]
\centering
\begin{tabular}{@{}lll@{}}
\toprule
\textbf{Library} & \textbf{Contents} & \textbf{Required For} \\
\midrule
Tc2\_Standard & TON, TOF, R\_TRIG, etc. & Basic PLC functions \\
Tc2\_System & ADSLOGSTR, system functions & Logging, diagnostics \\
Tc2\_Utilities & File I/O, advanced functions & Extended functionality \\
\bottomrule
\end{tabular}
\caption{Common TwinCAT Libraries}
\end{table}

\begin{tipbox}
To add a library reference, right-click on ``References'' in your PLC project, select ``Add library...'', and search for the library name. The Tc2\_Standard library should always be present for basic PLC functionality.
\end{tipbox}

% ============================================================
\section*{Quick Reference Card}
\addcontentsline{toc}{section}{Quick Reference Card}
% ============================================================

\subsection*{Timer Function Blocks}

\begin{lstlisting}
// TON - On Delay Timer
fbTON : TON := (PT := T#5S);
fbTON(IN := bTrigger);
// Q goes TRUE after 5 seconds if IN stays TRUE

// TOF - Off Delay Timer
fbTOF : TOF := (PT := T#5S);
fbTOF(IN := bTrigger);
// Q stays TRUE for 5 seconds after IN goes FALSE

// TP - Pulse Timer
fbTP : TP := (PT := T#5S);
fbTP(IN := bTrigger);
// Q goes TRUE for exactly 5 seconds on rising edge
\end{lstlisting}

\subsection*{Edge Detection}

\begin{lstlisting}
// Rising Edge
rtTrigger : R_TRIG;
rtTrigger(CLK := bInput);
IF rtTrigger.Q THEN
    // Executes once on FALSE->TRUE transition
END_IF

// Falling Edge
ftTrigger : F_TRIG;
ftTrigger(CLK := bInput);
IF ftTrigger.Q THEN
    // Executes once on TRUE->FALSE transition
END_IF
\end{lstlisting}

\subsection*{String Functions}

\begin{lstlisting}
// Concatenation
sResult := CONCAT(s1, s2);

// Substring
sLeft := LEFT(sSource, 5);     // First 5 chars
sRight := RIGHT(sSource, 3);   // Last 3 chars
sMid := MID(sSource, 2, 4);    // 4 chars from pos 2

// Length
nLen := LEN(sSource);
\end{lstlisting}

\subsection*{Time Literals}

\begin{table}[h]
\centering
\begin{tabular}{@{}ll@{}}
\toprule
\textbf{Literal} & \textbf{Value} \\
\midrule
T\#100MS & 100 milliseconds \\
T\#5S & 5 seconds \\
T\#2M30S & 2 minutes 30 seconds \\
T\#1H & 1 hour \\
T\#1D & 1 day \\
\bottomrule
\end{tabular}
\end{table}

\vfill

\begin{center}
\textit{Last Updated: \today}
\end{center}

\end{document}
